\section{System overview and design principles}
\label{sec:system}

QDCVR (v2.3) is a self-hosted knowledge-base management system for converting
heterogeneous documents into searchable, governable collections. A Python
parsing service, a TypeScript web application and an MCP tool server share
local storage and retrieval engines \cite{mcp}.

\subsection{Design principles: three enforceable contracts}
\label{sec:req}
\label{sec:principles}

Each contract names the mechanism that enforces it and the check that detects
a violation; enforcement is server-side throughout, so no contract depends on
an agent choosing to honour it.

\begin{description}
\item[Content governs placement.] \emph{Mechanism:} an ingest agent reads the
parsed content together with the library catalogue, assigns a parent or
sub-library, and raises an error when no scope applies. \emph{Check:}
descriptions and tags are re-read against the body, and three-source
consistency is verified after every lifecycle operation
(Sections~\ref{sec:acquisition} and~\ref{sec:organization}).
\item[Content governs admission.] \emph{Mechanism:} one lane-independent
judging tool returning a verdict with provenance, and lane policies that make
rank position advisory rather than authoritative. \emph{Check:} the gate fails
closed on missing text, tool errors and invalid verdicts, and the ablation of
Section~\ref{sec:e8} measures what removing it changes
(Sections~\ref{sec:lanes} and~\ref{sec:gate}).
\item[Absence is reported.] \emph{Mechanism:} an answer-layer contract in the
chat endpoint, not a similarity threshold. \emph{Check:} out-of-corpus probes
and unanswerable questions are graded as abstentions (Section~\ref{sec:e5}).
\end{description}

Four deployment requirements constrain all three contracts: expose failures
when parsing heterogeneous documents; keep document files and library records
consistent; require content evidence before a retrieved item supports an
answer; and remain operable on one workstation. Enforcement covers parse,
metadata and index post-conditions, the gate, and the lane tool policies.

\subsection{Architecture and subsystems}
\label{sec:architecture}

Fig.~\ref{fig:arch} shows the acquisition and query paths. Ingestion parses and
quality-checks documents; an agent uses the parsed content and library
catalogue to select a parent or sub-library; save and index operations then
update records and search indices. At query time, the chat endpoint supplies a
lane-specific tool policy, and the agent retrieves candidates and invokes the
gate before the answer layer returns cited evidence or not-found. The MCP
server exposes 95 registered tools. The three contracts of
Section~\ref{sec:principles} sit at different points along this path:
placement during ingestion, admission at the gate, the abstention report in
the answer layer.

The management path covers acquisition (Section~\ref{sec:acquisition}),
content-routed organization and lifecycle checks
(Section~\ref{sec:organization}), and the tool and runtime layer
(Section~\ref{sec:implementation}). The chat
endpoint defines pinned-search and all-library librarian policies. Hybrid is
a composed retrieval configuration, not a third built-in server policy; it
combines vector and catalog calls under an explicit directive. The document
graph supports navigation, not evidence, and the experience store uses the
same catalog and gate but was not evaluated at scale.

\begin{figure*}[t]
\centering
\includegraphics[width=\textwidth]{fig1_architecture}
\caption{QDCVR architecture. Ingestion parses and quality-checks documents,
and content-based library attribution precedes storage and indexing
(\textbf{C1}, content governs placement). At query time, the chat endpoint
applies a lane-specific tool policy; candidate retrieval is followed by
server-side content verification (\textbf{C2}, content governs admission),
after which the system returns cited evidence or an explicit not-found
response (\textbf{C3}, absence is reported). Authoritative records are
distinct from retrieval indices.}
\label{fig:arch}
\end{figure*}

\subsection{Data model and object lifecycle}
\label{sec:datamodel}
\label{sec:lifecycle}

Documents are stored as Markdown with extracted images. Three sources describe
library state: document files, the \texttt{.tree-fs.json} registry and each
library's \texttt{.knowledge-base.yml}. Ingestion proceeds through raw,
parsed, stored and indexed states, and is complete only after a content-derived
probe retrieves the document. Indexing is explicit and asynchronous; task
identifiers support polling, while post-operation audits detect stale indices
or disagreement among the three sources.
